\clearpage
\section{Методы естественно-языкового взаимодействия с базами данных}
\subsection{Постановка задачи интерпретации запроса}

Естественно-языковой интерфейс к базе данных (Natural Language Interface to Database, NLIDB) предназначен для преобразования высказывания пользователя в формальное представление, позволяющее получить требуемую информацию. В частном случае реляционной СУБД решается задача Text-to-SQL: по фразе на естественном языке и информации о схеме формируется SQL-запрос. В отличие от традиционного графического конструктора запросов, пользователь не обязан самостоятельно выбирать все операции реляционной алгебры. Однако за удобство такой формы взаимодействия приходится платить неоднозначностью исходного текста.

Например, высказывание «покажи активных клиентов с большими заказами» не задаёт однозначно смысл слов «активный» и «большой». Для получения корректного результата необходимо определить допустимые интерпретации этих понятий, выяснить период измерения активности, выбрать таблицы и условия агрегации. Нельзя считать верным преобразование, которое лишь порождает исполнимый SQL, не устанавливая соответствия между запросом и намерением пользователя.

Р.~В.~Посевкин и И.~А.~Бессмертный описывают диалоговый естественно-языковой интерфейс, преобразующий запросы в SQL \cite{posevkin2016}. В последующем исследовании Р.~В.~Посевкин использует семантическую модель базы данных и промежуточное представление для связывания фраз запроса с элементами схемы \cite{posevkin2018}. Эти работы значимы тем, что рассматривают не только синтаксис, но и механизм отображения смысловых конструкций на объекты данных. Для русскоязычных интерфейсов дополнительную сложность создают словоизменение, свободный порядок слов и вариативность предметной терминологии.

\subsection{Основные подходы к интерпретации запросов}

В литературе можно выделить три условных семейства подходов.\par\subsubsection{Подходы на основе правил и шаблонов}\par Первое основано на правилах и шаблонах. После лексического и синтаксического анализа из фразы выделяются сущности, атрибуты и операции; результат собирается с помощью заранее определённых конструкций. Преимуществами являются интерпретируемость и возможность ограничивать класс допустимых запросов, недостатками — высокая стоимость ручной настройки и слабая устойчивость к новым формулировкам.

\subsubsection{Семантические модели предметной области}\par Второе семейство опирается на семантические модели предметной области. Система использует связи между объектами, словарь терминов и сведения о структуре базы данных, чтобы сократить неоднозначность сопоставления. Работа Р.~В.~Посевкина показывает один из вариантов такого отображения \cite{posevkin2018}. В статье М.~Дж.~Аль-Мухаммеда и Д.~У.~Лонсдейла рассматривается адаптируемый агент для свободных естественно-языковых запросов, использующий онтологическое описание базы \cite{almuhammed2022}. Применение этого семейства подходов требует поддержки метаданных и правил соответствия, зато позволяет формализовать промежуточные решения.

\subsubsection{Обучаемые модели и Text-to-SQL}\par Третье семейство связано с обучаемыми моделями: классификаторами, средствами извлечения сущностей, трансформерами и большими языковыми моделями. И.~А.~Лапин и О.~Ю.~Сабинин рассматривают автоматическую генерацию модели базы данных из текстового задания с использованием распознавания сущностей и отношений \cite{lapin2024}. Эта задача не совпадает с Text-to-SQL, поскольку её результатом выступает модель данных, однако она демонстрирует применимость обработки естественного языка к восстановлению структуры предметной области. Т.~М.~Босенко исследует оптимизацию малых языковых моделей для генерации SQL-запросов, подчёркивая значение ограничений вычислительных ресурсов \cite{bosenko2025}.

Д.~А.~Евсеев описывает преобразование русскоязычных вопросов в запросы SPARQL с использованием синтаксического разбора и моделей для определения сущностей и отношений \cite{evseev2021}. SPARQL не является языком запросов к реляционным таблицам в смысле SQL, поэтому указанную работу следует рассматривать как методологически близкую, а не как эксперимент по SQL-генерации. Она полезна для обоснования многоэтапной интерпретации запроса: извлечения структуры высказывания, определения объектов и заполнения слотов формального запроса.

\subsection{Сопоставление подходов и их ограничений}

Ф.~Маджид и соавторы сравнивают естественно-языковые интерфейсы к базам данных и выделяют характеристики, существенные для пользователя, включая особенности ввода и качество получаемых запросов \cite{majeed2022}. Из сопоставления работ можно заключить, что универсального метода, одновременно гарантирующего точность произвольных запросов, прозрачность логики и отсутствие затрат на настройку предметной области, в рассмотренных публикациях не представлено. Это не утверждение о невозможности такого решения, а результат анализа выбранного корпуса литературы.

\begin{table}[htbp]
\caption{Сравнение основных способов интерпретации запросов (авторская систематизация)}
\label{tab:approaches}
\small
\begin{tabularx}{\textwidth}{|p{30mm}|X|X|}
\hline
\textbf{Подход} & \textbf{Сильные стороны} & \textbf{Ограничения} \\\hline
Шаблоны и правила & Контролируемая структура результата; объяснимость преобразования & Требуют ручной настройки; плохо масштабируются на вариативные формулировки \\\hline
Семантические модели & Учитывают значения предметных понятий и связи схемы & Нуждаются в актуальных словарях, онтологиях и метаданных \\\hline
Обучаемые модели & Лучше работают с разнообразным языком; позволяют переносить знания между формулировками & Возможны неверные ссылки на поля и ошибки смысловой интерпретации; требуется верификация \\\hline
Гибридный подход & Позволяет разделить генерацию и формальную проверку результата & Усложняет архитектуру и требует ясных правил согласования компонентов \\\hline
\end{tabularx}
\end{table}

Таблица~\ref{tab:approaches} представляет концептуальную классификацию, составленную для настоящего исследования на основе рассмотренных публикаций, а не численный рейтинг. Для содержательного сравнения в дальнейшем следует использовать единый набор запросов и схем, одинаковые ограничения на исполнение и раздельные метрики синтаксической и семантической корректности.

\subsection{Проверка намерения пользователя и безопасное исполнение}

Важным теоретическим различием является разграничение \emph{предложения операции} и её \emph{исполнения}. Пусть $u$ — высказывание пользователя, $S$ — описание схемы, $q$ — предложенная формальная операция, а $P$ — политика допустимости. Тогда концептуально процесс может быть представлен как
\begin{equation}
 q = G(u,S), \qquad a = V(q,S,P),
 \label{eq:query-flow}
\end{equation}
где $G$ — компонент генерации или интерпретации, $V$ — проверка предлагаемой операции, а $a$ — результат проверки: допустимо ли выполнение, требуется ли уточнение или подтверждение, следует ли отклонить запрос. Формула~\eqref{eq:query-flow} не задаёт конкретного алгоритма и используется только для разделения функций.

\begin{figure}[htbp]
\centering
\begin{tikzpicture}[
 box/.style={draw=black,rounded corners=2pt,minimum height=11mm,text width=30mm,align=center,font=\small},
 arr/.style={-{Stealth[length=2.5mm]},thick},node distance=5mm]
\node[box] (q) {Формулировка пользователя};
\node[box,right=of q] (g) {Интерпретация намерения};
\node[box,right=of g] (v) {Проверка и уточнение};
\node[box,right=of v] (e) {Исполнение операции};
\draw[arr] (q)--(g);
\draw[arr] (g)--(v);
\draw[arr] (v)--(e);
\node[draw,dashed,rectangle,align=center,font=\small,minimum width=36mm,minimum height=9mm,below=11mm of g] (m) {Метаданные и правила};
\draw[arr] (m.north)--(g.south);
\draw[arr] (m.east)-|(v.south);
\end{tikzpicture}
\caption{Концептуальные этапы обработки пользовательского запроса}
\label{fig:intent-process}
\end{figure}

Схема на рисунке~\ref{fig:intent-process} иллюстрирует разграничение ответственности компонентов. Она сформирована в рамках теоретического анализа и не представляет результаты измерений или уже реализованную архитектуру.

Такая схема помогает сформулировать проектный принцип: компонент, который предлагает SQL, не должен автоматически получать полномочия на произвольное выполнение. Минимальный уровень контроля включает анализ целевых таблиц и полей, определение типа операции, проверку допустимых действий и подтверждение изменений. Для запросов чтения также целесообразно ограничение размера выборки и времени выполнения; при ошибке интерпретации система должна уметь запросить уточнение вместо молчаливого выбора одного толкования.

Исследование О.~И.~Захаровой демонстрирует преимущества компонентной организации интеллектуальной обработки текстовых данных и конфигурируемых конвейеров \cite{zakharova2025}. В настоящей работе этот принцип рассматривается как архитектурный ориентир для разделения этапов интерпретации, проверки и доступа к данным. Само по себе применение компонентной архитектуры не доказывает безопасность: необходимы отдельная спецификация правил и проверка соответствия реализованного механизма этим правилам.

\subsection{Выводы по разделу}

Литература подтверждает существование нескольких научно обоснованных способов преобразования естественного языка в структурированные запросы. Для задачи SmartSQL наиболее содержательной представляется гипотеза о целесообразности комбинирования гибкой интерпретации с формальным контролем результата. Проверять эту гипотезу следует отдельно для запросов чтения и для операций изменения данных, поскольку последствия ошибок существенно различаются.